\section{Related works}\label{sec:related_works}

Acoustic signal analysis in combination with machine learning has attracted significant attention for fault detection in water-related infrastructure, as mechanical degradation frequently results in distinctive changes in emitted sound patterns. For water pumping systems, the authors in~\cite{sun2023pump} proposed a fault detection approach utilizing Mel Frequency Cepstral Coefficients (MFCC) and a Migration Learning Convolutional Neural Network (MLCNN). This method leverages both static MFCC features and their temporal derivatives to differentiate between normal and abnormal pump operating conditions, demonstrating the potential of acoustic features and deep learning even with limited labeled data. These results highlight the effectiveness of MFCC-based representations for pump condition monitoring, while maintaining compatibility with conventional computing resources.

Beyond pumping systems, recent research has investigated embedded and TinyML-based acoustic fault-detection models in related domains. The work of~\cite{lin2024real} introduced an application-specific embedded architecture that integrates spectrogram, Mel-filterbank energy (MFE), and MFCC feature extraction with Convolutional Neural Networks (CNN) for real-time defect identification in subsurface tiles. Similarly, authors of~\cite{atanane2023smart} examined TinyML-enabled leak detection in water distribution pipes, employing CNNs applied to scalogram images derived from vibration signals. These studies demonstrate the feasibility of conducting acoustic-based inference directly on low-power devices, thereby enabling local decision-making and reducing reliance on centralized processing infrastructures.

Although prior studies have established the suitability of acoustic features and embedded learning for fault detection, they have not systematically examined the trade-offs involved in deploying various audio preprocessing techniques under the strict memory and latency constraints of microcontroller-class devices. Specifically, comparative assessments of widely used features such as MFCC and MFE in fully on-device TinyML settings remain scarce. This gap motivates the current study, which investigates the end-to-end deployment of an acoustic-based fault detection model on a resource-constrained microcontroller, with an explicit evaluation of preprocessing strategies with respect to latency, memory footprint, and real-time deployability.
